%%%%%%%%%%%%%%%%%%%%%%%%%%%%%%%%%%%%%%%%%%%%%%%%%%%%%%%%%%%%%%%%%%%%%%%
% BAB 1
%%%%%%%%%%%%%%%%%%%%%%%%%%%%%%%%%%%%%%%%%%%%%%%%%%%%%%%%%%%%%%%%%%%%%%%
\mychapter{1}{BAB 1 PENDAHULUAN}

\section{Latar Belakang}

Wi-Fi telah menjadi teknologi komunikasi nirkabel yang digunakan secara luas di berbagai sektor. Saat ini, lebih dari 21 miliar perangkat telah mendukung teknologi Wi-Fi, dengan lebih dari 63\% lalu lintas internet global ditransmisikan melalui jaringan Wi-Fi \parencite{GERACI_WIFI_2025}. Kondisi tersebut menunjukkan bahwa kompetensi dalam merancang, mengimplementasikan, dan mengelola jaringan Wi-Fi menjadi semakin penting untuk memenuhi kebutuhan infrastruktur digital. Oleh karena itu, pengembangan kompetensi di bidang implementasi jaringan Wi-Fi \textit{enterprise} perlu didukung melalui kegiatan pelatihan berbasis praktik sehingga peserta tidak hanya memahami konsep jaringan, tetapi juga mampu melakukan konfigurasi, implementasi, dan validasi jaringan menggunakan perangkat yang umum digunakan di lingkungan industri.

Salah satu kegiatan yang mendukung pengembangan kompetensi tersebut adalah APIE (\textit{Asia Pacific Internet Engineering}) Pilot Project yang diselenggarakan dengan kolaborasi Direktorat Teknologi Informasi Universitas Brawijaya. Kegiatan ini merupakan program pelatihan berbasis \textit{hands-on} yang berfokus pada implementasi jaringan Wi-Fi \textit{enterprise} melalui serangkaian skenario praktikum. Dalam setiap skenario, peserta melakukan konfigurasi \textit{access point}, \textit{wireless LAN controller}, \textit{switch}, serta perangkat jaringan lainnya untuk membangun, mengonfigurasi, dan memverifikasi layanan jaringan sesuai dengan tujuan pembelajaran yang telah ditetapkan.

Untuk memastikan setiap skenario praktikum dapat digunakan sesuai dengan tujuan pembelajaran, lingkungan jaringan yang akan digunakan perlu dipersiapkan dan divalidasi sebelum pelaksanaan kegiatan. Proses tersebut mencakup perancangan jaringan, implementasi konfigurasi perangkat, pengujian konektivitas, serta verifikasi bahwa setiap skenario dapat dijalankan secara konsisten menggunakan perangkat yang sama dengan yang akan digunakan oleh peserta. Menurut \textcite{BRIDOVA_WIFI_DESIGN_2023}, hasil perancangan jaringan tetap memerlukan validasi melalui implementasi dan pengujian pada lingkungan nyata karena kondisi operasional dapat berbeda dengan hasil simulasi. Oleh karena itu, implementasi dan pengujian menjadi tahapan penting untuk memastikan lingkungan praktikum siap digunakan selama pelaksanaan APIE Pilot Project.

Berdasarkan uraian tersebut, Praktik Kerja Lapangan ini mengangkat judul \textit{"Implementasi dan Pengujian Skenario Jaringan Wi-Fi pada APIE Pilot Project di Universitas Brawijaya"}. Laporan ini membahas proses perancangan, implementasi, konfigurasi, pengujian, dan evaluasi skenario jaringan Wi-Fi sebagai bagian dari persiapan teknis APIE Pilot Project sehingga lingkungan praktikum yang digunakan telah memenuhi kebutuhan pembelajaran dan mendukung kelancaran pelaksanaan kegiatan.

\section{Rumusan Masalah}

Berdasarkan latar belakang yang telah diuraikan, rumusan masalah dalam Praktik Kerja 
Lapangan ini adalah sebagai berikut.

\begin{enumerate}[nosep]
\item Bagaimana proses perancangan dan implementasi skenario jaringan Wi-Fi pada APIE Pilot Project Universitas Brawijaya?
\item Bagaimana hasil pengujian dan verifikasi skenario jaringan Wi-Fi dalam memverifikasi bahwa setiap skenario telah berfungsi sesuai dengan kebutuhan pelatihan pada APIE Pilot Project Universitas Brawijaya?
\end{enumerate}


\section{Tujuan Penelitian}

Tujuan yang ingin dicapai dalam pelaksanaan Praktik Kerja Lapangan ini adalah sebagai berikut.

\begin{enumerate}[nosep]
\item Mendeskripsikan proses perancangan dan implementasi skenario jaringan Wi-Fi pada APIE Pilot Project Universitas Brawijaya.
\item Mendeskripsikan hasil pengujian dan verifikasi skenario jaringan Wi-Fi dalam memverifikasi bahwa setiap skenario telah berfungsi sesuai dengan kebutuhan pelatihan pada APIE Pilot Project Universitas Brawijaya.
\end{enumerate}

\section{Manfaat}
Manfaat yang diharapkan dari pelaksanaan Praktik Kerja Lapangan ini adalah sebagai berikut.

\subsection{Bagi Direktorat Teknologi Informasi Universitas Brawijaya}
Memberikan dokumentasi mengenai proses implementasi dan pengujian skenario jaringan Wi-Fi yang dapat dijadikan sebagai referensi dalam persiapan maupun penyelenggaraan kegiatan pelatihan berbasis jaringan pada masa mendatang.

\subsection{Bagi Penyelenggara APIE Pilot Project}
Memberikan dokumentasi hasil implementasi dan pengujian skenario jaringan Wi-Fi yang dapat dimanfaatkan sebagai bahan evaluasi dan pengembangan pelaksanaan APIE Pilot Project pada kegiatan berikutnya.

\subsection{Bagi Peserta}
Memberikan lingkungan praktikum jaringan Wi-Fi yang telah diimplementasikan dan diuji sehingga setiap skenario pelatihan dapat dijalankan sesuai dengan tujuan pembelajaran.

\subsection{Bagi Pembaca}
Memberikan referensi mengenai proses implementasi dan pengujian skenario jaringan Wi-Fi enterprise sebagai salah satu tahapan dalam mempersiapkan lingkungan pelatihan berbasis praktik.

\section{Batasan Masalah}

Supaya pembahasan lebih terarah, ruang lingkup Praktik Kerja Lapangan ini dibatasi pada hal-hal sebagai berikut.

\begin{enumerate}[nosep]
\item Objek pembahasan difokuskan pada implementasi dan pengujian skenario jaringan Wi-Fi yang digunakan pada APIE Pilot Project Universitas Brawijaya.
\item Implementasi yang dibahas meliputi perancangan topologi, konfigurasi perangkat jaringan, validasi konektivitas, serta pengujian setiap skenario sebelum digunakan dalam kegiatan pelatihan.
\item Pembahasan tidak mencakup penyusunan materi maupun modul pelatihan, melainkan hanya aspek teknis implementasi dan pengujian jaringan Wi-Fi.
\item Aktivitas selama pelaksanaan APIE Pilot Project dibahas sebatas dukungan teknis dan troubleshooting yang berkaitan dengan skenario jaringan Wi-Fi.
\end{enumerate}
